\documentclass[10pt,a4paper]{amsart}
\usepackage{amsmath,amssymb,amsthm,graphicx}
\usepackage{verbatim}
\newenvironment{reference}{\comment}{\endcomment}
%\newenvironment{reference}{}{}
\newenvironment{slogan}{\comment}{\endcomment}
\newenvironment{history}{\comment}{\endcomment}

% For commutative diagrams we use Xy-pic
\usepackage[all]{xy}

% We use 2cell for 2-commutative diagrams.
\xyoption{2cell}
% Theorem environments.
%
\theoremstyle{plain}
\newtheorem{theorem}[subsection]{Theorem}
\newtheorem{proposition}[subsection]{Proposition}
\newtheorem{lemma}[subsection]{Lemma}

\theoremstyle{definition}
\newtheorem{definition}[subsection]{Definition}
\newtheorem{example}[subsection]{Example}
\newtheorem{exercise}[subsection]{Exercise}
\newtheorem{situation}[subsection]{Situation}

\theoremstyle{remark}
\newtheorem{remark}[subsection]{Remark}
\newtheorem{remarks}[subsection]{Remarks}

\numberwithin{equation}{subsection}

% Macros
%
\def\lim{\mathop{\mathrm{lim}}\nolimits}
\def\colim{\mathop{\mathrm{colim}}\nolimits}
\def\Spec{\mathop{\mathrm{Spec}}}
\def\Hom{\mathop{\mathrm{Hom}}\nolimits}
\def\Ext{\mathop{\mathrm{Ext}}\nolimits}
\def\SheafHom{\mathop{\mathcal{H}\!\mathit{om}}\nolimits}
\def\SheafExt{\mathop{\mathcal{E}\!\mathit{xt}}\nolimits}
\def\Sch{\mathit{Sch}}
\def\Mor{\mathop{\mathrm{Mor}}\nolimits}
\def\Ob{\mathop{\mathrm{Ob}}\nolimits}
\def\Sh{\mathop{\mathit{Sh}}\nolimits}
\def\NL{\mathop{N\!L}\nolimits}
\def\CH{\mathop{\mathrm{CH}}\nolimits}
\def\proetale{{pro\text{-}\acute{e}tale}}
\def\etale{{\acute{e}tale}}
\def\QCoh{\mathit{QCoh}}
\def\Ker{\mathop{\mathrm{Ker}}}
\def\Im{\mathop{\mathrm{Im}}}
\def\Coker{\mathop{\mathrm{Coker}}}
\def\Coim{\mathop{\mathrm{Coim}}}

% Boxtimes
%
\DeclareMathSymbol{\boxtimes}{\mathbin}{AMSa}{"02}

%
% Macros for moduli stacks/spaces
%
\def\QCohstack{\mathcal{QC}\!\mathit{oh}}
\def\Cohstack{\mathcal{C}\!\mathit{oh}}
\def\Spacesstack{\mathcal{S}\!\mathit{paces}}
\def\Quotfunctor{\mathrm{Quot}}
\def\Hilbfunctor{\mathrm{Hilb}}
\def\Curvesstack{\mathcal{C}\!\mathit{urves}}
\def\Polarizedstack{\mathcal{P}\!\mathit{olarized}}
\def\Complexesstack{\mathcal{C}\!\mathit{omplexes}}
% \Pic is the operator that assigns to X its picard group, usage \Pic(X)
% \Picardstack_{X/B} denotes the Picard stack of X over B
% \Picardfunctor_{X/B} denotes the Picard functor of X over B
\def\Pic{\mathop{\mathrm{Pic}}\nolimits}
\def\Picardstack{\mathcal{P}\!\mathit{ic}}
\def\Picardfunctor{\mathrm{Pic}}
\def\Deformationcategory{\mathcal{D}\!\mathit{ef}}
\usepackage{cleveref}
\crefname{originallemma}{original lemma}{original lemmas}
\Crefname{originallemma}{Original lemma}{Original lemmas}
\makeatletter
\newlabel{lemma-twisted-inverse-image}{{0A9E}{source}{}{sourcepointer-1}{}}
\newlabel{lemma-twisted-inverse-image@cref}{{[originallemma][0][]0A9E}{[1][0][]source}{}{}{}}
\makeatother
\crefname{originallemma}{original lemma}{original lemmas}
\Crefname{originallemma}{Original lemma}{Original lemmas}
\makeatletter
\newlabel{perfect-lemma-affine-morphism-equivalence}{{0AVW}{source}{}{sourcepointer-2}{}}
\newlabel{perfect-lemma-affine-morphism-equivalence@cref}{{[originallemma][0][]0AVW}{[1][0][]source}{}{}{}}
\makeatother
\crefname{originallemma}{original lemma}{original lemmas}
\Crefname{originallemma}{Original lemma}{Original lemmas}
\makeatletter
\newlabel{more-morphisms-lemma-Noetherian-pseudo-coherent}{{0684}{source}{}{sourcepointer-3}{}}
\newlabel{more-morphisms-lemma-Noetherian-pseudo-coherent@cref}{{[originallemma][0][]0684}{[1][0][]source}{}{}{}}
\makeatother
\crefname{originallemma}{original lemma}{original lemmas}
\Crefname{originallemma}{Original lemma}{Original lemmas}
\makeatletter
\newlabel{more-morphisms-lemma-finite-pseudo-coherent}{{0AVX}{source}{}{sourcepointer-4}{}}
\newlabel{more-morphisms-lemma-finite-pseudo-coherent@cref}{{[originallemma][0][]0AVX}{[1][0][]source}{}{}{}}
\makeatother
\crefname{originallemma}{original lemma}{original lemmas}
\Crefname{originallemma}{Original lemma}{Original lemmas}
\makeatletter
\newlabel{perfect-lemma-quasi-coherence-internal-hom}{{0A6H}{source}{}{sourcepointer-5}{}}
\newlabel{perfect-lemma-quasi-coherence-internal-hom@cref}{{[originallemma][0][]0A6H}{[1][0][]source}{}{}{}}
\makeatother
\crefname{originallemma}{original lemma}{original lemmas}
\Crefname{originallemma}{Original lemma}{Original lemmas}
\makeatletter
\newlabel{lemma-twisted-inverse-image-bounded-below}{{0A9I}{source}{}{sourcepointer-6}{}}
\newlabel{lemma-twisted-inverse-image-bounded-below@cref}{{[originallemma][0][]0A9I}{[1][0][]source}{}{}{}}
\makeatother
\crefname{originallemma}{original lemma}{original lemmas}
\Crefname{originallemma}{Original lemma}{Original lemmas}
\makeatletter
\newlabel{lemma-sheafhom-affine-adjoint}{{0BV0}{source}{}{sourcepointer-7}{}}
\newlabel{lemma-sheafhom-affine-adjoint@cref}{{[originallemma][0][]0BV0}{[1][0][]source}{}{}{}}
\makeatother
\title{For a finite pseudo-coherent morphism the right adjoint of pushforward is the internal Hom out of the pushed forward structure sheaf}
\author{The Stacks Project Authors}
\date{stacks-project:duality.tex:local-snapshot-20261012; GNU FDL 1.2 or later}
\begin{document}
\maketitle
\noindent\textit{Editorial note.} All delivered mathematical text is the author's own, copied from the pinned Stacks Project chapter file duality.tex; no statement or proof step was rewritten, reordered or removed.\par\medskip
\noindent\textit{Editorial note.} Preamble provenance: the pinned chapter file opens with an input of the shared preamble and carries no preamble text of its own, so the project's single Stacks-wide preamble staged at sources/stacks-duality/preamble.tex is excerpted verbatim under the recipe fixed by LANE\_BRIEF 22.1, as three spans, lines 8-13 (verbatim package and the reference/slogan/history environments), lines 14-18 (Xy-pic setup) and lines 154-217 (the native theorem declarations and mathematical macros). Nothing was added, removed or redefined; the wrapper class and the amsmath/amssymb/amsthm/graphicx base packages are supplied by the assembler.\par\medskip
\noindent\textit{Editorial note.} The counted statement and proof contain no cross reference, no citation and no included file, so the delivered body is self contained.\par\medskip
\noindent\textit{Editorial note.} Canonical location of the counted proposition: Stacks tag 0AX2 (sources/stacks-duality/pinned-original-source.tex line 2763).\par\medskip
\noindent\textit{Allowed standard prerequisites (editorial).}\begin{itemize}
\item Pseudo-coherence of a finite morphism and of its pushforward of the structure sheaf. More on Morphisms, Lemmas labelled more-morphisms-lemma-Noetherian-pseudo-coherent and more-morphisms-lemma-finite-pseudo-coherent, cited at proof lines 2801-2809. Makes the internal Hom land in the required category.
\item Quasi-coherence of the internal Hom of a pseudo-coherent module. Derived Categories of Schemes, Lemma labelled perfect-lemma-quasi-coherence-internal-hom, cited at proof lines 2809-2813. Gives the mapping property.
\item The two adjunctions identifying the functors. Same chapter, Lemmas labelled lemma-twisted-inverse-image-bounded-below and lemma-sheafhom-affine-adjoint, cited at proof lines 2815-2819, and Derived Categories of Schemes, Lemma labelled perfect-lemma-affine-morphism-equivalence, cited in the statement. Shows both functors are right adjoint to the same functor.
\end{itemize}\medskip
\ifdefined\hypertarget\hypertarget{sourcepointer-1}{}\fi
\noindent\textit{Source reference.} Original lemma 0A9E: sources/stacks-duality/pinned-original-source.tex, source lines 653--653. Same-chapter pointer into the author environment lemma at source line 653; invoked at proof line None of the counted proof. It is retained only as a navigation pointer; every load-bearing step is inside the delivered span.\par\medskip
\ifdefined\hypertarget\hypertarget{sourcepointer-2}{}\fi
\noindent\textit{Source reference.} Original lemma 0AVW: sources/stacks-chow/tags.txt, source lines 13699--13699. Cross-chapter standard prerequisite located in the project tag list staged in the local snapshot; invoked at proof line None of the counted proof. It is retained only as a navigation pointer; every load-bearing step is inside the delivered span.\par\medskip
\ifdefined\hypertarget\hypertarget{sourcepointer-3}{}\fi
\noindent\textit{Source reference.} Original lemma 0684: sources/stacks-chow/tags.txt, source lines 7993--7993. Cross-chapter standard prerequisite located in the project tag list staged in the local snapshot; invoked at proof line 2784 of the counted proof. It is retained only as a navigation pointer; every load-bearing step is inside the delivered span.\par\medskip
\ifdefined\hypertarget\hypertarget{sourcepointer-4}{}\fi
\noindent\textit{Source reference.} Original lemma 0AVX: sources/stacks-chow/tags.txt, source lines 13700--13700. Cross-chapter standard prerequisite located in the project tag list staged in the local snapshot; invoked at proof line 2787 of the counted proof. It is retained only as a navigation pointer; every load-bearing step is inside the delivered span.\par\medskip
\ifdefined\hypertarget\hypertarget{sourcepointer-5}{}\fi
\noindent\textit{Source reference.} Original lemma 0A6H: sources/stacks-chow/tags.txt, source lines 12845--12845. Cross-chapter standard prerequisite located in the project tag list staged in the local snapshot; invoked at proof line 2792 of the counted proof. It is retained only as a navigation pointer; every load-bearing step is inside the delivered span.\par\medskip
\ifdefined\hypertarget\hypertarget{sourcepointer-6}{}\fi
\noindent\textit{Source reference.} Original lemma 0A9I: sources/stacks-duality/pinned-original-source.tex, source lines 726--726. Same-chapter pointer into the author environment lemma at source line 726; invoked at proof line 2797 of the counted proof. It is retained only as a navigation pointer; every load-bearing step is inside the delivered span.\par\medskip
\ifdefined\hypertarget\hypertarget{sourcepointer-7}{}\fi
\noindent\textit{Source reference.} Original lemma 0BV0: sources/stacks-duality/pinned-original-source.tex, source lines 2737--2737. Same-chapter pointer into the author environment lemma at source line 2737; invoked at proof line 2798 of the counted proof. It is retained only as a navigation pointer; every load-bearing step is inside the delivered span.\par\medskip
\begin{lemma}
\label{lemma-finite-twisted}
Let $f : Y \to X$ be a finite pseudo-coherent morphism of schemes
(a finite morphism of Noetherian schemes is pseudo-coherent).
The functor $R\SheafHom(f_*\mathcal{O}_Y, -)$ maps
$D_\QCoh^+(\mathcal{O}_X)$ into $D_\QCoh^+(f_*\mathcal{O}_Y)$.
If $X$ is quasi-compact and quasi-separated, then the diagram
$$
\xymatrix{
D_\QCoh^+(\mathcal{O}_X) \ar[rr]_a \ar[rd]_{R\SheafHom(f_*\mathcal{O}_Y, -)}
& & D_\QCoh^+(\mathcal{O}_Y) \ar[ld]^\Phi \\
& D_\QCoh^+(f_*\mathcal{O}_Y)
}
$$
is commutative, where $a$ is the right adjoint of
Lemma \ref{lemma-twisted-inverse-image} for $f$ and $\Phi$ is the equivalence
of Derived Categories of Schemes, Lemma
\ref{perfect-lemma-affine-morphism-equivalence}.
\end{lemma}

\begin{proof}
(The parenthetical remark follows from More on Morphisms, Lemma
\ref{more-morphisms-lemma-Noetherian-pseudo-coherent}.)
Since $f$ is pseudo-coherent, the $\mathcal{O}_X$-module $f_*\mathcal{O}_Y$
is pseudo-coherent, see More on Morphisms, Lemma
\ref{more-morphisms-lemma-finite-pseudo-coherent}.
Thus $R\SheafHom(f_*\mathcal{O}_Y, -)$ maps
$D_\QCoh^+(\mathcal{O}_X)$ into
$D_\QCoh^+(f_*\mathcal{O}_Y)$, see
Derived Categories of Schemes, Lemma
\ref{perfect-lemma-quasi-coherence-internal-hom}.
Then $\Phi \circ a$ and $R\SheafHom(f_*\mathcal{O}_Y, -)$
agree on $D_\QCoh^+(\mathcal{O}_X)$ because these functors are
both right adjoint to the restriction functor
$D_\QCoh^+(f_*\mathcal{O}_Y) \to D_\QCoh^+(\mathcal{O}_X)$. To see this
use Lemmas \ref{lemma-twisted-inverse-image-bounded-below} and
\ref{lemma-sheafhom-affine-adjoint}.
\end{proof}

\end{document}
